\chapter{สถาปัตยกรรมระบบตรวจวัดสุขภาพดิจิทัลและเซนเซอร์เอ็มอีเอ็มเอส}
\label{chap:ch01}

\section*{แผนบริหารการสอนประจำบทที่ 1}
\addcontentsline{toc}{section}{แผนบริหารการสอนประจำบทที่ 1}

\noindent\textbf{หัวข้อเนื้อหาประจำบท}
\begin{enumerate}
    \item ความสำคัญของระบบสารสนเทศสุขภาพส่วนบุคคลและการป้องกันโรคไม่ติดต่อเรื้อรัง
    \item หลักการฟิสิกส์และการทำงานของเซนเซอร์วัดความเร่งแบบไมโครอิเล็กโทรเมคานิกส์ (MEMS)
    \item สถาปัตยกรรมฮาร์ดแวร์และการเชื่อมต่อสัญญาณเซนเซอร์ในระบบสมาร์ทโฟน
    \item วงรอบการสุ่มสัญญาณตามทฤษฎีบทการชักตัวอย่างของไนควิสต์-แชนนอน
    \item สถาปัตยกรรมเลเยอร์ซอฟต์แวร์และการทำงานร่วมกันระหว่าง Flutter Platform Channel กับระบบปฏิบัติการ
    \item มาตรฐานความปลอดภัยและความเป็นส่วนตัวของข้อมูลสุขภาพส่วนบุคคล
\end{enumerate}

\noindent\textbf{วัตถุประสงค์เชิงพฤติกรรม}
\begin{enumerate}
    \item อธิบายหลักการทางฟิสิกส์ของโครงสร้างหวีตัวเก็บประจุภายในเซนเซอร์ MEMS Accelerometer ได้อย่างถูกต้อง
    \item วิเคราะห์บทบาทของความถี่การสุ่มตัวอย่างต่อความถูกต้องในการตรวจจับการเคลื่อนไหวของมนุษย์ได้อย่างถูกต้อง
    \item จำแนกลำดับชั้นสถาปัตยกรรมซอฟต์แวร์ตั้งแต่ระดับฮาร์ดแวร์จนถึงแอปพลิเคชันระดับผู้ใช้ได้อย่างครบถ้วน
    \item ประยุกต์ใช้แนวทางการจัดเก็บข้อมูลสุขภาพส่วนบุคคลให้สอดคล้องกับมาตรฐานความปลอดภัยระดับสากลได้อย่างถูกต้อง
\end{enumerate}

\noindent\textbf{กิจกรรมการเรียนการสอน}
\begin{enumerate}
    \item การบรรยายเชิงทฤษฎีประกอบการสาธิตภาพจำลองโครงสร้างหวีตัวเก็บประจุจุลภาคของเซนเซอร์ MEMS
    \item การอภิปรายกลุ่มเกี่ยวกับผลกระทบของพฤติกรรมเนือยนิ่งต่อสุขภาพและแนวทางการใช้เทคโนโลยีสมาร์ทโฟนเข้าช่วยตรวจติดตาม
    \item การทดลองอ่านค่าสัญญาณความเร่งสามแกนสดจากสมาร์ทโฟนผ่านโปรแกรมตรวจสอบสัญญาณ
    \item การสรุปบทเรียนและร่วมทำแบบฝึกหัดทบทวนท้ายบท
\end{enumerate}

\noindent\textbf{สื่อการเรียนการสอน}
\begin{enumerate}
    \item เอกสารคำสอนและสไลด์บรรยายเรื่องสถาปัตยกรรมเซนเซอร์สุขภาพดิจิทัล
    \item สมาร์ทโฟนระบบปฏิบัติการแอนดรอยด์ที่มีเซนเซอร์วัดความเร่งสามแกน
    \item แพลตฟอร์มจำลองสัญญาณความเร่งและซอฟต์แวร์บันทึกข้อมูลแบบเรียลไทม์
    \item ภาพจำลองโครงสร้างจุลภาคเสมือนจริงของชิปเซนเซอร์ซิลิคอน
\end{enumerate}

\noindent\textbf{การประเมินผล}
\begin{enumerate}
    \item การประเมินผลการมีส่วนร่วมในการอภิปรายและการตอบคำถามในชั้นเรียน
    \item การประเมินผลสัมฤทธิ์จากการทำใบงานการวิเคราะห์คุณลักษณะเซนเซอร์ MEMS
    \item การตรวจความถูกต้องของแบบฝึกหัดคำถามทบทวนท้ายบทที่ 1
\end{enumerate}

\newpage

\section{ความสำคัญของระบบสารสนเทศสุขภาพส่วนบุคคล}
ในศตวรรษที่ 21 ภัยคุกคามด้านสาธารณสุขที่สำคัญที่สุดประการหนึ่งคือกลุ่มโรคไม่ติดต่อเรื้อรัง (Non-Communicable Diseases หรือ NCDs) เช่น โรคเบาหวานชนิดที่ 2 โรคความดันโลหิตสูง โรคหลอดเลือดสมองและหัวใจ ตลอดจนภาวะอ้วนลงพุง ซึ่งมีสาเหตุหลักมาจากการใช้ชีวิตประจำวันที่มีพฤติกรรมเนือยนิ่ง (Sedentary Lifestyle) เป็นเวลานาน ขาดกิจกรรมทางกายที่เพียงพอตามคำแนะนำขององค์การอนามัยโลก (World Health Organization: WHO)

การติดตามพฤติกรรมสุขภาพแบบดั้งเดิมมักอาศัยการจดบันทึกด้วยตนเอง (Self-Reporting) ซึ่งมักมีความลำเอียงในการรายงาน (Recall Bias) สูงและไม่สะท้อนความจริงอย่างต่อเนื่อง การนำเทคโนโลยีสารสนเทศสุขภาพดิจิทัล (Digital Health Informatics) เข้ามาประยุกต์ใช้ โดยอาศัยสมาร์ทโฟนที่บุคคลพกพาติดตัวอยู่ตลอดเวลา จึงเป็นแนวทางที่มีประสิทธิภาพสูงสุดในการเปลี่ยนผ่านจากการรักษาเมื่อเจ็บป่วย (Curative Care) ไปสู่การเฝ้าระวังเชิงป้องกันก่อนเกิดโรค (Preventive Healthcare)

\section{หลักการฟิสิกส์ของเซนเซอร์วัดความเร่งแบบ MEMS}
สมาร์ทโฟนยุคใหม่ได้รับการติดตั้งระบบเครื่องกลไฟฟ้าจุลภาค (Micro-Electro-Mechanical Systems: MEMS) ซึ่งรวมเอาชิ้นส่วนกลไกขนาดไมครอนเข้ากับวงจรอิเล็กทรอนิกส์บนแผ่นเวเฟอร์ซิลิคอนชิปเดียวกัน หัวใจสำคัญในการตรวจจับการก้าวเดินคือ เซนเซอร์วัดความเร่งแบบสามแกน (3-Axis MEMS Accelerometer)

\begin{theorybox}
\textbf{หลักการทำงานของโครงสร้างหวีตัวเก็บประจุ (Capacitive Comb-Drive)}

ภายในชิป MEMS มีมวลทดสอบ (Proof Mass, $m$) ถูกแขวนลอยอยู่ด้วยสปริงซิลิคอนจุลภาค (Micro-Spring, ค่าคงที่สปริง $k$) และมีฟันหวีซิลิคอนตัวนำไฟฟ้าที่ประกบสลับกันระหว่างส่วนที่เคลื่อนที่ได้กับส่วนที่ยึดแน่นกับตัวเรือนชิป

เมื่อร่างกายเกิดการก้าวเดิน เกิดความเร่ง $a$ ตามกฎข้อที่สองของนิวตัน $F = ma$ แรงเฉื่อยจะผลักให้มวลทดสอบขยับออกจากตำแหน่งสมดุลเป็นระยะกระจัด $x = \frac{ma}{k}$ ส่งผลให้ระยะห่างระหว่างฟันหวีเปลี่ยนแปลงไป ก่อให้เกิดการเปลี่ยนแปลงค่าความจุไฟฟ้า ($\Delta C$) ดังความสัมพันธ์
\begin{equation}
C = \frac{\varepsilon_0 \varepsilon_r A}{d} \quad \implies \quad \Delta C \approx \varepsilon_0 \varepsilon_r A \left( \frac{1}{d - x} - \frac{1}{d + x} \right) \approx \frac{2 \varepsilon_0 \varepsilon_r A}{d^2} x
\end{equation}
โดยที่ $\varepsilon_0$ คือค่าสภาพยอมของสุญญากาศ $\varepsilon_r$ คือค่าคงที่ไดอิเล็กทริก $A$ คือพื้นที่ผิวหน้าฟันหวีที่สบกัน และ $d$ คือระยะห่างเริ่มต้นระหว่างแผ่นฟันหวี วงจรแปลงประจุจะแปลงการเปลี่ยนแปลงค่าความจุไฟฟ้านี้ให้กลายเป็นสัญญาณแรงดันไฟฟ้าอนาล็อก จากนั้นจะถูกแปลงเป็นสัญญาณตัวเลขดิจิทัล (ADC) ในพิกัดสามมิติ $a_x, a_y, a_z$
\end{theorybox}

\section{วงรอบการสุ่มสัญญาณและทฤษฎีบทไนควิสต์-แชนนอน}
การเคลื่อนไหวก้าวเดินและวิ่งของมนุษย์ปกติมีความถี่ของการสวิงแขนและขยับก้าวขาอยู่ในย่านความถี่ตั้งแต่ $0.5\text{ Hz}$ จนถึงไม่เกิน $3.5\text{ Hz}$ และการกระแทกของส้นเท้า (Heel Strike) อาจก่อให้เกิดฮาร์มอนิกความถี่สูงไม่เกิน $15\text{ Hz}$

ตามทฤษฎีบทการสุ่มสัญญาณของไนควิสต์-แชนนอน (Nyquist-Shannon Sampling Theorem) อัตราความถี่ในการสุ่มสัญญาณดิจิทัล ($f_s$) จะต้องมีค่าอย่างน้อยสองเท่าของความถี่สูงสุดที่มีนัยสำคัญในสัญญาณ ($f_{\max}$) เพื่อป้องกันปัญหาภาพซ้อนสัญญาณความถี่เท็จ (Aliasing)
\begin{equation}
f_s \ge 2 f_{\max} \quad \implies \quad f_s \ge 2 \times 15\text{ Hz} = 30\text{ Hz}
\end{equation}
ในสถาปัตยกรรมของแอปพลิเคชัน \textbf{Smart Health Tracker} ผู้วิจัยได้กำหนดอัตราการสุ่มตัวอย่างไว้ที่ $50\text{ Hz}$ (คาบการสุ่ม $T_s = 20\text{ ms}$) ซึ่งมีความถี่สูงกว่าขีดจำกัดไนควิสต์ถึง $1.67$ เท่า ส่งผลให้สามารถจับยอดพีคของแรงกระแทกจากส้นเท้าได้อย่างแม่นยำ โดยไม่กินพลังงานแบตเตอรี่ของสมาร์ทโฟนมากเกินไป

\begin{table}[H]
\centering
\caption{การเปรียบเทียบอัตราการสุ่มสัญญาณเซนเซอร์ความเร่งกับประสิทธิภาพการใช้พลังงาน}
\label{tab:sensor_sampling_rates}
\small
\begin{tabularx}{\textwidth}{|c|c|X|c|}
\toprule
\textbf{ความถี่สุ่ม ($f_s$)} & \textbf{คาบเวลา ($T_s$)} & \textbf{ความเหมาะสมในการตรวจจับพฤติกรรมสุขภาพ} & \textbf{การบริโภคพลังงาน} \\
\midrule
$10\text{ Hz}$ & $100\text{ ms}$ & สุ่มช้าเกินไป ยอดพีคตกหล่น เกิดความผิดพลาดสูง & ต่ำมาก ($<0.5\text{ mA}$) \\
$20\text{ Hz}$ & $50\text{ ms}$ & พอใช้สำหรับการเดินช้า แต่อาจพลาดจังหวะการวิ่งเร็ว & ต่ำ ($0.8\text{ mA}$) \\
$\mathbf{50\text{ Hz}}$ & $\mathbf{20\text{ ms}}$ & \textbf{จุดสมดุลสูงสุด ตรวจจับพีคคมชัด แม่นยำ ครอบคลุมการวิ่ง} & \textbf{ปานกลาง ($1.5\text{ mA}$)} \\
$100\text{ Hz}$ & $10\text{ ms}$ & สัญญาณละเอียดมากเกินความจำเป็น สัญญาณรบกวนความถี่สูงมาก & สูง ($3.2\text{ mA}$) \\
\bottomrule
\end{tabularx}
\end{table}

\section{สถาปัตยกรรมซอฟต์แวร์และการเชื่อมต่อแพลตฟอร์ม}
สถาปัตยกรรมของระบบ Smart Health Tracker พัฒนาบนเฟรมเวิร์ก Flutter โดยใช้ภาษา Dart ออกแบบโครงสร้างแบบแยกชั้น (Layered Clean Architecture) เพื่อให้การเชื่อมต่อสัญญาณฮาร์ดแวร์มีเสถียรภาพและไม่กระทบต่อประสิทธิภาพการแสดงผลบนหน้าจอ (UI Thread)

\begin{enumerate}
    \item \textbf{Hardware and Kernel Driver Layer} ระดับชิปเซนเซอร์และไดรเวอร์เคอร์เนลลินุกซ์ของระบบปฏิบัติการแอนดรอยด์ ทำหน้าที่อ่านค่าแรงดันไฟฟ้าและปรับเทียบค่าความโน้มถ่วงศูนย์ ($0g$ Offset)
    \item \textbf{Android Sensor HAL and Java/Kotlin Layer} เลเยอร์ฮาร์ดแวร์แอบสแตรกชันของระบบปฏิบัติการแอนดรอยด์ โดยเรียกใช้คลาส \texttt{SensorManager} และลงทะเบียนรับฟังเหตุการณ์ผ่าน \texttt{SensorEventListener} ด้วยตัวเลือกความเร็ว \texttt{SENSOR\_DELAY\_GAME} หรือ \texttt{SENSOR\_DELAY\_UI}
    \item \textbf{Flutter Platform Channel Layer} กลไกการส่งผ่านข้อมูลสัญญาณความเร่งแบบอะซิงโครนัสข้ามสภาพแวดล้อมภาษา โดยใช้แพ็กเกจ \texttt{sensors\_plus} แปลงเป็นสตรีมข้อมูล \texttt{accelerometerEventStream()} เข้าสู่ฝั่ง Dart
    \item \textbf{Digital Signal Processing Engine} คอร์ประมวลผลสัญญาณ คำนวณเวกเตอร์รวมยูคลิด กรองความถี่ต่ำ กรองยอดพีคแบบพลวัต และตรวจสอบระยะพักฟื้น
    \item \textbf{Exercise Science and Health State} วิเคราะห์และแปลงข้อมูลก้าวเป็นพลังงานกิโลแคลอรี ระยะทางก้าว อัตราความถี่ก้าว และสถานะการนั่งนิ่ง
    \item \textbf{Adaptive Bento Grid UI Layer} ส่วนติดต่อผู้ใช้ที่แสดงผลข้อมูลเป็นบล็อกการ์ดที่สวยงาม ปรับขนาดตามหน้าจอโดยอัตโนมัติ
\end{enumerate}

\begin{figure}[H]
\centering
\begin{tikzpicture}[scale=0.88, every node/.style={transform shape}]
  % Architecture blocks
  \node[draw=rbruNavy, fill=rbruNavy!10, rounded corners=4pt, line width=1pt, minimum width=13cm, minimum height=1cm, align=center] (app) at (0, 5) {\textbf{Presentation Layer: Adaptive Bento Grid UI \& Kinetic Neon Avatar}};
  \node[draw=rbruCyan, fill=rbruCyan!10, rounded corners=4pt, line width=1pt, minimum width=13cm, minimum height=1cm, align=center] (biz) at (0, 3.6) {\textbf{Domain Layer: Exercise Science, Mifflin-St Jeor BMR, Cadence \& Hydration}};
  \node[draw=rbruEmerald, fill=rbruEmerald!10, rounded corners=4pt, line width=1pt, minimum width=13cm, minimum height=1cm, align=center] (dsp) at (0, 2.2) {\textbf{Physics \& DSP Engine: 3D Norm, Dynamic Peak Filter \& Refractory Guard}};
  \node[draw=rbruSlate, fill=rbruSlate!10, rounded corners=4pt, line width=1pt, minimum width=13cm, minimum height=1cm, align=center] (hal) at (0, 0.8) {\textbf{Platform Bridge: Flutter MethodChannel \& sensors\_plus Stream (50 Hz)}};
  \node[draw=black, fill=black!10, rounded corners=4pt, line width=1pt, minimum width=13cm, minimum height=1cm, align=center] (hw) at (0, -0.6) {\textbf{Hardware Layer: 3-Axis MEMS Comb-Drive Silicon Accelerometer}};

  % Arrows
  \draw[->, >=stealth, line width=1.5pt, color=rbruNavy] (hw) -- (hal);
  \draw[->, >=stealth, line width=1.5pt, color=rbruSlate] (hal) -- (dsp);
  \draw[->, >=stealth, line width=1.5pt, color=rbruEmerald] (dsp) -- (biz);
  \draw[->, >=stealth, line width=1.5pt, color=rbruCyan] (biz) -- (app);
\end{tikzpicture}
\caption{สถาปัตยกรรมระบบหลายชั้นของแอปพลิเคชัน Smart Health Tracker}
\label{fig:app_system_architecture}
\end{figure}

\section{ความมั่นคงปลอดภัยและความเป็นส่วนตัวของข้อมูลสุขภาพ}
ข้อมูลเกี่ยวกับพฤติกรรมการเคลื่อนไหวและตำแหน่งการก้าวเดินจัดเป็นข้อมูลสุขภาพส่วนบุคคลที่มีความอ่อนไหว (Sensitive Personal Data) ภายใต้พระราชบัญญัติคุ้มครองข้อมูลส่วนบุคคล พ.ศ. 2562 (PDPA) สถาปัตยกรรมของ Smart Health Tracker จึงยึดหลักการประมวลผล ณ อุปกรณ์ปลายทาง (On-Device Local Processing)

ข้อมูลดิบจากเซนเซอร์ความเร่งจะถูกคำนวณและสกัดคุณลักษณะภายในหน่วยความจำแรมของสมาร์ทโฟนทันที และผลรวมสถิติรายวันจะถูกจัดเก็บลงในพื้นที่จัดเก็บข้อมูลภายในแบบเข้ารหัสของแอปพลิเคชัน (Private Encrypted SharedPreferences) โดยไม่มีการส่งข้อมูลดิบออกไปยังเซิร์ฟเวอร์ภายนอก ช่วยป้องกันการดักจับข้อมูลและการระบุตัวตนของผู้ใช้ได้อย่างสมบูรณ์

\section*{คำถามทบทวนท้ายบทที่ 1}
\addcontentsline{toc}{section}{คำถามทบทวนท้ายบทที่ 1}
\begin{enumerate}
    \item จงอธิบายการเปลี่ยนแปลงของโครงสร้างหวีตัวเก็บประจุในเซนเซอร์ MEMS เมื่อร่างกายมีการเคลื่อนไหวพร้อมเขียนสมการความสัมพันธ์
    \item เพราะเหตุใดความถี่ในการสุ่มสัญญาณที่ $50\text{ Hz}$ จึงถือเป็นจุดสมดุลที่เหมาะสมในการตรวจจับการก้าวเดินเมื่อเทียบกับเกณฑ์ไนควิสต์
    \item จงอธิบายลำดับการส่งผ่านข้อมูลจากเซนเซอร์ฮาร์ดแวร์จนปรากฏเป็นตัวเลขบนหน้าจอแอปพลิเคชัน Flutter
    \item เพราะเหตุใดการประมวลผลข้อมูลสัญญาณบนอุปกรณ์ของผู้ใช้ (On-Device Processing) จึงสอดคล้องกับมาตรฐานความปลอดภัยของข้อมูลสุขภาพส่วนบุคคล
\end{enumerate}
\clearpage
